% !TEX TS-program = lualatexmk
% !TEX parameter =  --shell-escape

\documentclass[vespers_book_la_eng.tex]{subfiles}

\ifcsname preamble@file\endcsname
  \setcounter{page}{\getpagerefnumber{M-vespers_book_la_eng_epiphany_octave}}
\fi

\begin{document}
\festum{epiphany of the lord}
\addcontentsline{toc}{chapter}{Epiphany of the Lord and its Octave}\phantomsection
%\addcontentsline{toc}{section}{Epiphany of the Lord}\phantomsection

\doubleclass{I}[with a Privileged Octave of the II Order]

\atvespers{i}

\lastplace{II}{116}

\psalmus{116}{116_7c2}{116_7}{116_en}

\chapterlateng
\scripture{Isaiæ 60, 1.}

\texteparacol{\lettrine{S}{u}rge, illumináre Jerúsalem, quia venit lumen tuum,~* et glória Dómini super te orta est.}{english}{\noindent Arise, be enlightened, O Jerusalem: for thy light is come, and the glory of the Lord is risen upon thee.}

\label{hy_crudelis_herodes_deum_solesmes_1961} \phantomsection
\hymnlateng

\gscore[]{3.}{hy_crudelis_herodes_deum_solesmes_1961}
\textes{crudelis_herodes_en.tex}

\anothertone

\gscore[]{8.}{hy_crudelis_herodes_deum_another_chant_solesmes_1961}

\textes{octave_versiculus_rubrique}

\label{versiculus_vesperis_epiphianiae_domini}\phantomsection
   \smallscore{versiculus_vesperis_epiphianiae_domini}
   
   \begin{enpars}\noindent \vv The kings of Tarshish and of the isles shall bring presents.

\noindent \rr  The kings of Arabia and Saba shall offer gifts.\end{enpars}

\gscore[Ad Magnif.]{Ant. 8. G}{an_magi_videntes_stellam_solesmes_1961}

\antiphona{english}{When the wise men saw the star, they said one unto another: This is the sign of the great King; let us go and search diligently for him, and present unto him gifts, gold, and frankincense, and myrrh. Alleluia}

\label{oratio_epiphaniae} \phantomsection

\collectlateng

\texteparacol{
\lettrine{D}{e}us, qui hodiérna die Unigénitum tuum géntibus stella duce revelásti:~† concéde propítius; ut qui jam te ex fide cognóvimus,~* usque ad contemplándam spéciem tuæ celsitúdinis perducámur. Per eúmdem Dóminum.}{english}{\noindent 
O God, who by the leading of a star didst, as on this day, manifest thine Only-Begotten Son to the Gentiles, mercifully grant that we, which know thee now by faith, may after this life have the fruition of thy glorious godhead:
Through the same Lord.}

\atvespers{ii}
\firstantiphon{2. D}
\initialscore{an_ante_luciferum_genitus_solesmes_1961}
\antiphona{english}{The Lord our Saviour, begotten before the day-star, and before the ages, is this day made manifest in the world}
\psalmus{109}{109_2}{109_2}{109_en}

 \gscore[2. Ant.]{1. g2}{an_venit_lumen_tuum_solesmes_1961}
 \antiphona{english}{O Jerusalem, thy light is come, and the glory of the Lord is risen upon thee, and the Gentiles shall walk in thy light. Alleluia}
\psalmus{110}{110_1g2}{110_1}{110_en}

  \gscore[3. Ant.]{1. g2}{an_apertis_thesauris_solesmes_1961}
  \antiphona{english}{When the wise men had opened their treasures, they presented unto the Lord gold, frankincense, and myrrh. Alleluia}
 \psalmus{111}{111_1g2}{111_1}{111_en}

\gscore[4. Ant.]{4. E}{an_maria_et_flumina_solesmes_1961}
\antiphona{english}{O ye seas and floods, bless ye the Lord. O ye wells, bless ye the Lord. Alleluia}
\psalmus{112}{112_4E}{112_4E}{112_en}

 \gscore[5. Ant.]{7. c2}{an_stella_ista_solesmes_1961}
\antiphona{english}{This star shineth as a flame of fire, and pointeth out God, the King of kings: the Wise Men beheld it, and offered gifts unto the great King}
\psalmus{113}{113_7c2}{113_7}{113_en}

\caphymnvven{I}

\gscore[Ad Magnif.]{Ant. 1. D}{an_tribus_miraculis_solesmes_1961}

\antiphona{english}{This day we keep a holiday in honour of three wonders, this day a star led the wise men to the manger; this day at the marriage, water was made wine; this day was Christ, for our salvation, pleased to be baptized of John in Jordan. Alleluia}

\rubrique{Collect \normaltext{Deus qui hodiérna,} as above, \normaltext{\pageref{oratio_epiphaniae}.}}

\litdateen{Sunday within the octave the Epiphany}

\bigtitle{Holy Family of Jesus, Mary, and Joseph.}
\addcontentsline{toc}{section}{Holy Family of Jesus, Mary, and Joseph}\phantomsection
\doublemajor

\rubrique{When the Octave of Epiphany falls on a Sunday, on the preceding Saturday the office of the Holy Family is said, and on Friday I Vespers of the Holy Family are said, with the commemoration of the preceding day within the octave and of the Sunday within the octave, as below.}

\rubrique{However, if a feast which is double of the I class occurs on this Saturday, the office of the Holy Family with the commemoration of that Sunday is anticipated to the nearest Friday on which the office of the octave would otherwise have been said; and in the office of both the feast double of the I class and of the Holy Family, a commemoration of the current day within the octave is made.}

\rubrique{Finally, from January 7 through 12 inclusive, when a Sunday and a feast which is double of the I class occur at the same time, at both Vespers of the feast, the commemoration is first of the Holy Family, whose Office. with all and individual rights has been substituted for the office of the Sunday in perpetuity, then of the Sunday and the octave.}

\newpage
\atvespersminor{I}

\firstantiphon{1. g}
\initialscore{an_jacob_autem_genuit_solesmes_1961}
\antiphona{english}{And Jacob begot Joseph the husband of Mary, of whom was born Jesus, who is called Christ}
\psalmus{109}{109_1g}{109_1}{109_en}

 \gscore[2. Ant.]{7. c}{an_angelus_domini_apparuit_holy_family_solesmes_1961}
\antiphona{english}{The angel of the Lord appeared to Joseph in his sleep, saying: Joseph, son of David, fear not to take unto thee Mary thy wife, for that which is conceived in her, is of the Holy Ghost}
\psalmus{112}{112_7c}{112_7}{110_en}

  \gscore[3. Ant.]{7. d}{an_pastores_venerunt_solesmes_1961}
 \antiphona{english}{The shepherds came with haste; and they found Mary and Joseph, and the infant lying in the manger}
  \psalmus{121}{121_7d}{121_7}{121_en}

  \gscore[4. Ant.]{2. D}{an_magi_intrantes_solesmes_1961}
\antiphona{english}{The Magi, entering into the house, found the child with Mary his mother}
\label{ps_126_2}
  \psalmus{126}{126_2}{126_2}{126_en}

 \gscore[5. Ant.]{1. f}{an_erat_pater_ejus_solesmes_1961}
\antiphona{english}{His father and mother were wondering at those things which were spoken concerning him}
\psalmus{147}{147_1f}{147_1}{147_en}

\chapterlateng \label{cap_san_fam} \phantomsection
\scripture{Luc. 2, 51.}

\texteparacol{\lettrine{D}{e}scéndit Jesus cum María et Joseph, et venit Názareth,~* et erat súbditus illis.}{english}{\noindent Jesus went down with Mary and Joseph, and came to Nazareth, and was subject to them.}

\hymnlateng \label{hy_o_lux_beata_caelitum_solesmes_1961} \phantomsection
\gscore[]{2.}{hy_o_lux_beata_caelitum_solesmes_1961}
\textes{o_lux_beata_caelitum_en}

\smallscore{versiculus_1_vespers_holy_family}

\begin{enpars}\noindent \vv Blessed are they that dwell in thy house, O Lord.

\noindent \rr They shall praise thee for ever and ever.
 \end{enpars}
 
 \gscore[Ad Magnif.]{Ant. 8. G}{an_verbum_caro_factum_est_solesmes_1961}
\antiphona{english}{The Word was made flesh, and dwelt among us, full of grace and truth; and of his fulness we all have received, and grace for grace, alleluia}

\label{oratio_san_fam} \phantomsection

\collectlateng

 \texteparacol{
\lettrine{D}{o}mine Jesu Christe, qui Maríæ et Joseph súbditus, domésticam vitam ineffabílibus virtútibus consecrásti:~† fac nos, utriúsque auxílio, Famíliæ sanctæ tuæ exémplis ínstrui;~* et consórtium cónsequi sempitérnum. Qui vivis et regnas cum Deo Patre.}{english}{\noindent
Lord Jesus Christ, who subject to Mary and Joseph, didst consecrate family life by thy unspeakable virtues, aid us by their united intercession to profit by the examples of thy Holy Family, and attain to their everlasting companionship. Who livest and reignest with God the Father.}

\rubrique{Commemoration of the octave of the Epiphany, as below, \normaltext{\pageref{octave_epiphaniae}.}}

\rubrique{Then one is made of the I Sunday after the Epiphany.}

\gscore[Ant.]{2.}{an_remansit_puer_solesmes_1961}
\antiphona{english}{The child Jesus remained in Jerusalem; and his parents knew it not. And thinking that he was in the company, they came a day's journey, and sought him among their kinsfolks and acquaintance}

\label{omnes_de_saba}
\texteparacol{\vv  Omnes de Saba vénient, allelúia.

\rr Aurum et thus deferéntes, allelúia.}{english}{\noindent \vv All they from Saba shall come. Alleluia.

\noindent \rr They shall bring gold and incense. Alleluia.}

\collectlateng
%[findent=0.05em]

\label{oratio_dom_1_pe}
\texteparacol{\lettrine{V}{o}ta quǽsumus Dómine, supplicántis pópuli cælésti pietáte proséquere:~† ut et quæ agénda sunt, vídeant,~* et ad implénda quæ víderint, convaléscant.
Per Dóminum.}{english}{\noindent Attend, O Lord, we beseech thee, of thy heavenly mercy, to the desires of thy suppliant people; and grant that they may both perceive what they ought to do, and may have strength to fulfill the same. Through our Lord.}
 
\atvespersminor{II}

\rubrique{When at I Vespers only a commemoration of the feast was made, the antiphons and psalms are those of I Vespers.}

\gscore[1. Ant.]{8. G}{an_post_triduum_solesmes_1961}
\antiphona{english}{After three days, they found him in the temple, sitting in the midst of the doctors, hearing them, and asking them questions}
\psalmus{109}{109_8G}{109_8}{109_en}

\gscore[2. Ant.]{4. E}{an_dixit_mater_jesu_ad_illum_solesmes_1961}
%%actually needs to go under Vespers of the Epiphany
\antiphona{english}{The mother of Jesus said to him: Son, why hast thou done so to us? Behold thy father and I have sought thee sorrowing}
\psalmus{112}{112_4E}{112_4E}{112_en}

%\rubrique{Ps. \normaltext{Laudáte Púeri, \pageref{M-112_4E}.}}

%%does it need the prefix if it's in the same subfile. TBD!

\gscore[3. Ant.]{8. G}{an_descendit_jesus_cum_eis_solesmes_1961}
\antiphona{english}{Jesus went down with them, and came to Nazareth, and was subject to them}
\psalmus{121}{121_8G}{121_8}{121_en}

\gscore[4. Ant.]{2. D}{an_et_jesus_proficiebat_solesmes_1961}
\antiphona{english}{And Jesus advanced in wisdom, and age, and grace with God and men}
\rubrique{Ps. \normaltext{Nisi Dóminus, \pageref{ps_126_2}.}}

\gscore[5. Ant.]{8. G}{an_et_dicebant_solesmes_1961}
\begin{enpars}\noindent Ant.  And they said: How came this man by this wisdom and miracles? Is not this the carpenter's son?\end{enpars}
\psalmus{147}{147_8G}{147_8}{147_en}

\rubrique{Chapter \normaltext{Descéndit,} as at I Vespers, \normaltext{\pageref{hy_o_lux_beata_caelitum_solesmes_1961}.}}

\rubrique{Hymn \normaltext{O lux beáta Cǽlitum, \pageref{hy_o_lux_beata_caelitum_solesmes_1961}.}}

\smallscore{versiculus_2_vespers_holy_family}

\begin{enpars}\noindent \vv All thy children shall be taught of the Lord.

\noindent \rr And great shall be the peace of thy children.
 \end{enpars}
 
  \gscore[Ad Magnif.]{Ant. 8. G}{an_maria_autem_solesmes_1961}
\antiphona{english}{But Mary kept all these words, pondering them in her heart}
 
\rubrique{Collect \normaltext{Dómine Jesu Christe,} as above, \normaltext{\pageref{oratio_san_fam}.}}

\rubrique{Commemoration of the octave of the Epiphany, as below, \normaltext{\pageref{octave_epiphaniae}.}}

\rubrique{Then one is made of the Sunday.}

\gscore[Ant.]{8.}{an_fili_quid_fecisti_solesmes_1961}
\begin{enpars}\noindent Ant. Son, why hast thou thus dealt with us? Behold, thy father and I have sought thee sorrowing. How is it that ye sought me? Wist ye not that I must be about my Father's business?\end{enpars}

\rubrique{\normaltext{℣. Omnes de Saba, \pageref{omnes_de_saba}.} Collect Vota, \pageref{oratio_dom_1_pe}.}

\rubrique{¶ If the Octave Day of the Epiphany occurs on Sunday, nothing is celebrated of the Sunday on the octave day itself, but at I Vespers only, a commemoration for II Vespers of Sunday is made, if on January 12 was celebrated the office of the Holy Family as above.}

\addcontentsline{toc}{section}{Days within the Octave of the Epiphany} \label{octave_epiphaniae}\phantomsection
\bigtitle{January 7. Second day within the Octave.}
 
\gscore[Ad Magnif.]{Ant. 7. a}{an_videntes_stellam_solesmes_1961}
 \antiphona{english}{When the wise men saw the star, they rejoiced with exceeding great joy; and when they were come into the house, they presented unto the Lord gold, frankincense, and myrrh. Alleluia}
 
 \rubrique{When the day within the octave is commemorated in the office of the Holy Family:}
 
 \vv  Reges Tharsis et ínsulæ múnera ófferent.

\rr Reges Arabum et Saba dona addúcent.

\rubrique{Collect \normaltext{Deus qui hodiérna,} as above, \normaltext{\pageref{oratio_epiphaniae}.}}
 
\bigtitle{January 8. Third day.}

\gscore[Ad Magnif.]{Ant. 8. c}{an_lux_de_luce_solesmes_1961}
 \antiphona{english}{O Christ, thou Light of Light, thou art made manifest, and the wise men have presented unto thee gifts. Alleluia, Alleluia, Alleluia}
 
\bigtitle{January 9. Fourth day.}

\gscore[Ad Magnif.]{Ant. 8. G}{an_interrogabat_magos_solesmes_1961}
 \antiphona{english}{Herod inquired of the wise men: What is the sign that ye have seen that a King is born? We have seen the shining of the star, the glory whereof enlighteneth the whole world}
 
\bigtitle{January 10. Fifth day.}

\gscore[Ad Magnif.]{Ant. 8. c}{an_omnes_de_saba_solesmes_1961}
 \antiphona{english}{All they from Saba shall come, they shall bring gold and frankincense, alleluia, alleluia}
 
 \rubrique{Commemoration of \normaltext{ S. Hygini,} Pope et Martyr. Ant. \normaltext{Iste sanctus, ℣. Glória et honóre.}}
 
%\rubrique{Commemoratio \normaltext{ S. Hygini,} Papæ et Martyris. Ant. \normaltext{Iste sanctus, \pageref{M-an_cs_iste_sanctus_pro_lege_solesmes_1961}, ℣. Glória et honóre.}}

%%this is in the antiphonal as is. page for St Telesphorus refers (properly) to the commemorations section.

%\gscore[Ant.]{8.}{an_cs_iste_sanctus_pro_lege_solesmes_1961}
%
%\vv Glória et honóre coronásti eum Dómine.
%
%\rr Et constituísti eum super ópera mánuum tuárum.

\collectlateng

\texteparacol{\lettrine{I}{n}firmitátem nostram réspice, omnípotens Deus:~† et quia pondus própriæ actiónis gravat,~* beáti Hygíni Mártyris tui atque Pontíficis intercéssio gloriósa nos prótegat.
Per Dóminum.}{english}{\noindent Mercifully consider our weakness, o almighty God, and whereas by the burden of our sins we are sore, let Hyginus thy martyr be mercifully pleased to deliver us from all things which may hurt our bodies, and from all evil thoughts which may defile our souls.
Through our Lord.}

\bigtitle{January 11. Sixth day.}

\gscore[Ad Magnif.]{Ant. 1. f}{an_admoniti_magi_solesmes_1961}
 \antiphona{english}{The wise men being warned in a dream, departed into their own country another way}
 
\bigtitle{January 12. Seventh day.}

\rubrique{Vespers of the following Octave Day.}

\bigtitle{Octave of the Epiphany.}
\addcontentsline{toc}{section}{Octave of the Epiphany}\phantomsection

\rubrique{All as on the feast, except the collect.}

\collectlateng

\texteparacol{
\lettrine{D}{e}us, cujus Unigénitus in substántia nostræ carnis appáruit:~† præsta, quǽsumus; ut per eum, quem símilem nobis foris agnóvimus,~* intus reformári mereámur. 
Qui tecum vivit et regnat in unitáte.}{english}{\noindent
O God, whose only-begotten Son appeared in the substance of our flesh, grant, we beseech thee, that we who acknowledge his outward likeness to us may deserve to be inwardly refashioned in his image. Who with thee liveth and reigneth, in the unity.}

\end{document}